\documentclass[11pt]{article}
\usepackage[letterpaper,margin=1in]{geometry}
\usepackage[T1]{fontenc}
\usepackage{amsmath,amssymb,amsthm,lmodern,microtype}
\usepackage[hidelinks]{hyperref}

\newtheorem{theorem}{Theorem}
\newtheorem{lemma}[theorem]{Lemma}
\newtheorem{proposition}[theorem]{Proposition}
\newtheorem{corollary}[theorem]{Corollary}
\newtheorem{remark}[theorem]{Remark}

\title{Spanning Packings in Small Tuza Counterexamples}
\author{}
\date{}

\begin{document}
\maketitle

\begin{abstract}
Let $\nu(G)$ denote the maximum number of edge-disjoint triangles in a finite simple graph $G$, and let $\tau(G)$ denote the minimum number of edges meeting every triangle of $G$. A counterexample to Tuza's conjecture is \emph{minimal} if $\tau(G)>2\nu(G)$ and every proper subgraph satisfies $\tau\leq 2\nu$. We prove that in every minimal counterexample with $\nu(G)\leq 9$, a maximum triangle packing maximizing vertex support spans $G$, and the union of every maximum triangle packing is connected. In particular, $|V(G)|\leq 2\nu(G)+1$. Combined with the $\nu\leq 7$ theorem and Kandasamy's maximum-degree-ten theorem, this restricts the first unresolved case $\nu(G)=8$ to graph orders $12$ through $17$, and $\nu(G)=9$ to order at most $19$. We also establish a general upper bound $|V(G)|\leq 2\nu(G)+\lfloor(\nu(G)-1)/2\rfloor$ whenever a maximum-support packing fails to span $G$.
\end{abstract}

For a finite simple graph $G$, let $\nu(G)$ be the maximum number of edge-disjoint triangles and let $\tau(G)$ be the minimum number of edges meeting every triangle. Call $G$ a \emph{minimal counterexample} if $\tau(G)>2\nu(G)$ and every proper subgraph satisfies $\tau\leq 2\nu$.

\begin{theorem}\label{thm:main}
In any minimal counterexample $G$ with $\nu(G)\leq 9$, a maximum triangle packing covering the most vertices spans $G$, and the union of its triangles is connected. In particular, $|V(G)|\leq 2\nu(G)+1$. The first unresolved packing number is $\nu(G)=8$, for which only orders $12$ through $17$ need to be searched, while $\nu(G)=9$ has order at most $19$.
\end{theorem}

\section{Reductions and small links}

Let $G$ be a minimal counterexample to Tuza's conjecture. Minimality rules out isolated vertices and edges lying in fewer than two triangles: deleting an isolated vertex or a triangle-free edge leaves $\nu$ and $\tau$ unchanged, and deleting the three edges of a unique triangle containing an edge reduces $\nu$ by at least one and $\tau$ by at most two (the other two edges of the triangle cover all lost triangles). In particular, every link $L_v=G[N(v)]$ has minimum degree at least two.

A nonempty matching $M$ in a link component with vertex set $W\subseteq N(v)$ and a vertex set $Q\subseteq W$ of size at most $|M|$ covering all edges of the component outside $M$ give a reduction: removing the spokes from $v$ to $W$ and the edges of $M$ produces a proper subgraph whose packings extend by the $|M|$ triangles of $M$ and whose covers extend by $M$ together with the spokes from $v$ to $Q$, contradicting minimality.

\begin{lemma}\label{lem:min-nine-edges}
Every link $L_v$ of a minimal counterexample $G$ has at least nine edges, and a link with exactly nine edges is necessarily $K_5$ minus one edge.
\end{lemma}

\begin{proof}
No link component can be a cycle: a maximum alternating matching on an even cycle leaves a matching, and on an odd cycle leaves a two-edge path and isolated edges; both residuals have vertex cover number at most the matching size. Components of order four other than $C_4$ have a two-edge matching whose deletion leaves edges covered by the endpoints of one matching edge.

By Puleo's reductions (Lemmas~2.6 and~2.7 of~\cite{Puleo}), $\delta(G)\geq 5$, and for links of orders five or six the complement has maximum degree at most one and edge count different from two. Thus a five-vertex link has at least nine edges ($K_5$ or $K_5-e$), and a six-vertex link has at least twelve edges.

For links of orders seven, eight, or nine, $\delta(L_v)\geq 2$ gives $|E(L_v)|\geq |V(L_v)|$. Equality would mean $L_v$ is $2$-regular (a union of cycles), which is reducible. The only remaining cases with at most nine edges have $(|V(L_v)|,|E(L_v)|)\in\{(7,8),(7,9),(8,9)\}$. The exact link checker~\cite{SmallLinks} enumerates every labeled graph of each size and verifies a whole-link matching $M$ and residual vertex cover $Q$ with $|Q|\leq |M|$ on all $11{,}550$, $59{,}990$, and $140{,}420$ labeled graphs, respectively, of minimum degree at least two. Finally, any link of order at least ten has at least ten edges by $\delta(L_v)\geq 2$. Therefore every link has at least nine edges, with equality only for $K_5-e$.
\end{proof}

\section{Every maximum packing has connected union}

\begin{proposition}\label{prop:connected}
In any minimal counterexample $G$, every maximum triangle packing $\mathcal{P}$ has connected union and covers at most $2\nu(G)+1$ vertices.
\end{proposition}

\begin{proof}
Let $\mathcal{P}$ be a maximum triangle packing. Every triangle of $G$ uses at least one edge of $\mathcal{P}$, or it could enlarge $\mathcal{P}$. Form an auxiliary graph $\Gamma$ on $\mathcal{P}$, joining two packed triangles whenever some triangle of $G$ contains an edge from each. Every triangle of $G$ then meets packed edges from exactly one connected component of $\Gamma$.

For a component of $\Gamma$ containing $s$ packed triangles, let $G_s$ be obtained from $G$ by deleting the edges of all packed triangles outside this component. The triangles of $G_s$ are precisely the triangles of $G$ meeting packed edges in this component, and $\nu(G_s)=s$ (since any larger packing in $G_s$ together with the deleted packed triangles would enlarge $\mathcal{P}$). If $\Gamma$ had at least two components, each $G_s$ would be a proper subgraph of $G$ and hence would satisfy $\tau(G_s)\leq 2s$; the union of their covers would cover $G$ at cost at most $2\nu(G)$, a contradiction.

Thus $\Gamma$ is connected. Adjacent triangles in $\Gamma$ share a vertex (the common endpoint of their two edges in the witnessing triangle of $G$), so the union of the triangles of $\mathcal{P}$ is connected. Adding the $\nu(G)$ triangles one at a time along a spanning tree of $\Gamma$ adds at most two new vertices per step after the first triangle, yielding at most $3+2(\nu(G)-1)=2\nu(G)+1$ vertices.
\end{proof}

\section{Spanning maximum packings for packing number at most nine}

Let $\mathcal{P}$ be a maximum triangle packing of $G$ maximizing the vertex support $|V(\bigcup\mathcal{P})|$. Call a vertex in $V(\bigcup\mathcal{P})$ \emph{private} if it lies in exactly one triangle of $\mathcal{P}$, and \emph{repeated} if it lies in at least two. Let $X=V(G)\setminus V(\bigcup\mathcal{P})$ be the set of uncovered vertices.

\begin{lemma}\label{lem:uncovered-links}
The set $X$ is independent. For every $x\in X$, every edge of $L_x$ is a packed edge joining two repeated vertices whose packed triangle has a private third vertex. In particular, each packed triangle of $\mathcal{P}$ contributes at most one edge to $\bigcup_{x\in X}E(L_x)$.
\end{lemma}

\begin{proof}
For $x\in X$, no edge incident with $x$ is packed, so every triangle through $x$ uses a packed edge opposite $x$. Since every edge of $G$ lies in a triangle, $X$ is independent and every edge of $L_x$ is a packed edge.

If $uv\in E(L_x)$ lies in the packed triangle $uvw\in\mathcal{P}$, then $w$ must be private; otherwise replacing $uvw$ with $xuv$ would yield a maximum packing with strictly larger vertex support. Now suppose some private vertex $v$ were adjacent to $x$. Since $\delta(L_x)\geq 2$ and every edge of $L_x$ is packed, both neighbors of $v$ in its unique packed triangle lie in $L_x$, so all three edges of that triangle lie in $L_x$. Applying the replacement argument to all three edges shows that all three vertices are private, so none has any other packed edge; hence they form a triangle component of $L_x$, contradicting minimality.

Consequently, every vertex of $L_x$ is repeated, and every edge $uv\in E(L_x)$ lies in a packed triangle $uvw\in\mathcal{P}$ with $w$ private. Since $w$ cannot lie in any $L_{x'}$ ($x'\in X$), the triangle $uvw$ contributes at most the single edge $uv$ to $\bigcup_{x\in X}E(L_x)$.
\end{proof}

\begin{proof}[Proof of Theorem~\ref{thm:main}]
Let $\mathcal{P}$ be a maximum-support maximum packing of $G$, with uncovered set $X$. If $X\neq\varnothing$, pick $x\in X$. By Lemmas~\ref{lem:min-nine-edges} and~\ref{lem:uncovered-links}, $\nu(G)\geq |E(L_x)|\geq 9$, so $X=\varnothing$ whenever $\nu(G)\leq 8$.

Now suppose $\nu(G)=9$ and $X\neq\varnothing$. Then $|E(L_x)|=9$ for every $x\in X$, so by Lemma~\ref{lem:min-nine-edges} all nine packed triangles have two repeated vertices and one private vertex, and their nine opposite edges induce $F\cong K_5-e$ on a common set $R$ of five repeated vertices, with $N(x)=R$ and $G[R]=F$. Write each packed triangle as $p_e+e$ for $e\in E(F)$, where the nine private vertices $p_e$ are distinct.

We claim there is no edge between distinct private vertices $p_e$ and $p_f$. Suppose $p_ep_f\in E(G)$, and consider any triangle $p_ep_fw$. Since no uncovered vertex has a private neighbor (Lemma~\ref{lem:uncovered-links}) and three private vertices cannot span a packed edge, we must have $w=r\in R$. If $r$ belongs to neither $e$ nor $f$, then $p_ep_fr$ avoids all packed edges. If $r$ belongs to exactly one of them, say $f$, then $p_ep_fr$ uses a packed edge only from $p_f+f$; replacing $p_f+f$ by $p_ep_fr$ and $x+f$ yields ten edge-disjoint triangles, a contradiction. Hence $r$ must be a common endpoint of $e$ and $f$. Because distinct edges share at most one vertex, $p_ep_f$ lies in at most one triangle, contradicting minimality.

By $\delta(G)\geq 5$, each private vertex is therefore adjacent to all five vertices of $R$ and has no other neighbors. Fix one private vertex $p$ and one $x\in X$. Label $R=\{0,1,2,3,4\}$ so that $01$ is the missing edge of $F=G[R]$. Take the four edge-disjoint triangles $x02$, $x13$, $p03$, and $p12$. Consider the seven-edge set
\[
C=\{02,03,12,13,23,x4,p4\}.
\]
Delete $C$ and all remaining spokes at $x$ and $p$ to obtain a proper subgraph $H$. All four chosen triangles lie in the deleted edges, so $\nu(G)\geq\nu(H)+4$. Since $N(x)=N(p)=R$ and $01\notin E(G)$, every triangle meeting $\{x,p\}$ either contains $4$ (covered by $x4$ or $p4$) or uses an edge of $G[\{0,1,2,3\}]$ (covered by $\{02,03,12,13,23\}$). Thus $C$ covers all lost triangles, giving $\tau(G)\leq\tau(H)+7<2\nu(G)$, a contradiction. Hence $X=\varnothing$ for $\nu(G)=9$ as well.

Combining $X=\varnothing$ with Proposition~\ref{prop:connected} gives $|V(G)|\leq 2\nu(G)+1$. By the all-order theorem for $\nu\leq 7$~\cite{PackingSeven}, the first unresolved packing number is $\nu(G)=8$. By Kandasamy's maximum-degree-ten theorem~\cite{Kandasamy}, $\Delta(G)\geq 11$ and hence $|V(G)|\geq 12$. Thus a minimal counterexample with $\nu(G)=8$ has order between $12$ and $17$, and one with $\nu(G)=9$ has order at most $19$.
\end{proof}

\begin{remark}
At order $17$ with $\nu(G)=8$, the bipartite incidence graph between the $17$ vertices and the $8$ packed triangles of $\mathcal{P}$ is connected with $25$ nodes and $24$ edges, hence a tree. Removing the node for a packed triangle $uvz\in\mathcal{P}$ leaves three branches. If $z$ is repeated and $uvw$ is a triangle of $G$ using packed edges only from $uvz$, then $w$ must lie in the branch containing $z$: otherwise replacing $uvz$ by $uvw$ would leave the $z$-branch nonempty (as $z$ is repeated) and disconnected from the rest of the packing, contradicting Proposition~\ref{prop:connected}.
\end{remark}

\section{Further constraints on uncovered vertices}

\begin{proposition}
In any minimal counterexample $G$, for every nonempty independent set $W\subseteq V(G)$, the union $B(W)=\bigcup_{w\in W}E(L_w)$ satisfies $|B(W)|\geq 2|W|+1$. Consequently, for any packing number $\nu=\nu(G)$, if the uncovered set $X$ of a maximum-support packing is nonempty, then $|X|\leq\lfloor(\nu-1)/2\rfloor$ and $|V(G)|\leq 2\nu+\lfloor(\nu-1)/2\rfloor$.
\end{proposition}

\begin{proof}
Suppose $|B(W)|\leq 2|W|$ for some nonempty independent set $W$, and choose $W$ inclusion-minimal. Every proper nonempty subset $U\subsetneq W$ satisfies $|B(U)|\geq 2|U|+1\geq |U|$. For $W$ itself, if $|W|=1$ then $L_w$ is nonempty, while if $|W|>1$ removing one vertex gives $|B(W)|\geq 2|W|-1\geq |W|$. By Hall's theorem, there are distinct edges $e_w\in E(L_w)$ for $w\in W$. Because $W$ is independent, the $|W|$ triangles $w+e_w$ are pairwise edge-disjoint. Deleting $B(W)$ and all spokes at $W$ yields a proper subgraph $H$ with $\nu(G)\geq\nu(H)+|W|$, while $B(W)$ alone covers all lost triangles at cost $|B(W)|\leq 2|W|$, contradicting minimality.

Now suppose $X\neq\varnothing$ for a maximum-support packing $\mathcal{P}$. By Lemma~\ref{lem:uncovered-links}, $X$ is independent and $|B(X)|\leq\nu$, so $2|X|+1\leq\nu$ and $|X|\leq\lfloor(\nu-1)/2\rfloor$. Pick $x\in X$; since $\delta(L_x)\geq 2$, the link $L_x$ contains a cycle. Because each edge of that cycle comes from a distinct packed triangle of $\mathcal{P}$, the cycle lifts to a cycle in the connected bipartite incidence graph of $\mathcal{P}$, which has $3\nu$ edges and $|V(\bigcup\mathcal{P})|+\nu$ nodes. Hence $|V(\bigcup\mathcal{P})|+\nu\leq 3\nu$, i.e., $|V(\bigcup\mathcal{P})|\leq 2\nu$, yielding $|V(G)|\leq 2\nu+\lfloor(\nu-1)/2\rfloor$.
\end{proof}

\section{Computational replay}

Run \texttt{python3 check\_small\_links.py} with Python~3.10 or later. The standard-library checker enumerates all labeled graphs of orders and sizes $(7,8)$, $(7,9)$, and $(8,9)$, and verifies a whole-link matching $M$ and residual vertex cover $Q$ with $|Q|\leq |M|$ on all $11{,}550$, $59{,}990$, and $140{,}420$ labeled graphs, respectively, of minimum degree at least two~\cite{SmallLinks}.

\begin{thebibliography}{9}
\bibitem{Kandasamy}
A.~Kandasamy, \emph{Tuza's conjecture for graphs of maximum degree at most 10},
\href{https://arxiv.org/abs/2610.15269}{arXiv:2610.15269}, 2026.

\bibitem{Puleo}
G.~J.~Puleo, \emph{Tuza's conjecture for graphs with maximum average degree less than 7},
\href{https://doi.org/10.1016/j.ejc.2015.03.006}{European J. Combin. 49 (2015), 134--152}.

\bibitem{PackingSeven}
\emph{Tuza's inequality for triangle packing number at most seven}, companion manuscript,
\href{run:packing-seven.pdf}{packing-seven.pdf}.

\bibitem{SmallLinks}
Exact small-link reduction checker,
\href{https://github.com/claude-make-no-mistakes/tuza-counterexample-search/blob/main/check_small_links.py}{\texttt{check\_small\_links.py}}, and finite verification log,
\href{https://github.com/claude-make-no-mistakes/tuza-counterexample-search/blob/main/small-links.jsonl}{\texttt{small-links.jsonl}}.
\end{thebibliography}

\end{document}
